\subsubsection{The two square-zero six-dimensional exceptions}

\begin{lemma}\label{B:six-squarezero-target-modules}
Let $A$ be one of the following radical-square-zero rational
algebras, with orthogonal primitive idempotents $e_1,e_2,e_3$.
Suppose $\dim W=5$, $\dim V=2$, and
\[
 \bigl(\dim\operatorname{Tor}_0^A(W,V),
       \dim\operatorname{Tor}_1^A(W,V),
       \dim\operatorname{Tor}_2^A(W,V)\bigr)=(2,0,1).
\]
The conclusions below concern every such module pair, without
assuming that the displayed algebra basis is a monoid basis.
\begin{enumerate}
\item For the triangle algebra, take
$a\in e_1Ae_2$, $b\in e_2Ae_3$, $c\in e_1Ae_3$ and let $V$
be the two-dimensional module on which $c$ is the nonzero arrow.
Either every balanced form has rank at most one, or $W$ is
faithful over $A/(a)$ or over $A/(b)$.  In the first faithful
case the vertex dimensions of $W$ are $(1,1,3)$; in the second
they are $(2,1,2)$.
\item For the parallel-chain algebra, take
$a\in e_1Ae_2$ and $b,c\in e_2Ae_3$, with $b,c$ independent.
Let $V$ have $b$ acting as the identity between its one-dimensional
vertex spaces and $c$ acting as zero.  Either every balanced form
has rank at most one, or $W$ is faithful over $A/(a)$, with vertex
dimensions $(1,1,3)$.
\end{enumerate}
Every quotient algebra appearing in the faithful alternatives is
hereditary, as is each of its further algebra quotients.
\end{lemma}
\begin{proof}
Write $W_i=We_i$.  In the triangle case the exact projective
resolution of $V$ is
\[
 0\longrightarrow Ae_1\xrightarrow{a}Ae_2
 \xrightarrow{b}Ae_3\longrightarrow V\longrightarrow0.
\]
Here the first arrow sends $e_1$ to $a$, the second sends $e_2$
to $b$, and the last map retains $e_3,c$ and kills $b$.
After tensoring, the complex is
$W_1\xrightarrow{a}W_2\xrightarrow{b}W_3$.
Put $r=\operatorname{rank}a$ and $s=\operatorname{rank}b$.
The three Tor dimensions imply
\[
 \dim W_1=r+1,\qquad \dim W_2=r+s,\qquad
 \dim W_3=s+2.
\]
Their sum is five, so $r+s=1$.

A balanced form is specified by a functional on
$W_3/\operatorname{im}b$.  Its value on the $e_1$ coordinate of
$V$ is obtained by composing that functional with
$c:W_1\to W_3$.  If $c=0$, all forms therefore factor through
the same one-dimensional quotient of $V$ and have rank at most
one.  Otherwise, if $(r,s)=(0,1)$, the action kills $a$ and is
faithful over $A/(a)$: each of its three vertex spaces is nonzero,
and its remaining two arrows $b,c$ act nontrivially in different
Peirce matrix blocks.  If $(r,s)=(1,0)$, the identical argument
gives faithfulness over $A/(b)$.  The displayed vertex dimensions
follow from the three equations above.

In the parallel-chain case, the projective resolution has the
same form with its second arrow $c$ in place of $b$.
Put $r=\operatorname{rank}a$ and $s=\operatorname{rank}c$;
again $r+s=1$ with the same vertex dimensions.  If $r=1$, the
map $a:W_1\to W_2$ is onto, and the relation $ab=0$ forces
$b:W_2\to W_3$ to be zero.  Every balanced form then factors
through the top coordinate of $V$.  If $r=0$, the source of
both $b,c:W_2\to W_3$ has dimension one, and $c$ is nonzero.
When their images are dependent, $b$ is zero after passage to
$W_3/\operatorname{im}c$, so all forms again have the same
one-dimensional row factor.  When they are independent, the
three idempotent actions and the two arrow actions are linearly
independent, and $W$ is faithful over $A/(a)$.

The triangle quotients are the path algebras of two-arrow forks.
The parallel-chain quotient is a two-parallel-arrow algebra
together with an isolated vertex.  These quivers have no
composable pair of arrows.  Every quotient merely deletes
vertices or takes a quotient of an arrow space and is again
a path algebra of this form; hence it is hereditary.
\end{proof}

\begin{lemma}[The literal arrow-ideal gate]
\label{B:six-squarezero-literal-ideals}
Among the $458$ actual row-module cases in the two square-zero
exceptions of Proposition~\ref{B:six-row-census}, the faithful
alternatives of Lemma~\ref{B:six-squarezero-target-modules}
give $636$ one-dimensional arrow ideals to test.  All $280$
parallel-chain ideals and $348$ of the $356$ triangle ideals
are killed by the faithful transformation-pair image of every
monoid realizing the corresponding module pair.  Their images
therefore have zero rational fourth homology and all degree-two
products in the original monoids are zero.

The only eight exceptions to this literal-ideal test have the
source indices and ideals listed in the next proposition.
\end{lemma}
\begin{proof}
The checker \path{code/d_six_square_zero_ideals.py} uses the
complete zero-seven census, the exact primitive idempotents,
and the actual row representations in
\path{code/outputs/b_three_point_row_syzygies.json}.
If $t$ is the top vertex of the row module, $m$ is the vertex
of its one-dimensional first syzygy, and $s$ is the unique
vertex with $e_sJe_m\ne0$, the first tested ideal is
$I=e_sJe_m$.  In a triangle it also tests $I=e_mJe_t$.
Each is checked by exact multiplication to be a one-dimensional
two-sided ideal.

Here is the complete logical content of each successful finite
certificate.  The ideal is spanned by distinguished basis vectors
lying in it and differences $p-q$ of distinguished basis vectors
with the same image modulo it.  A basis vector in $I$ has zero
augmentation operator on both $V$ and $W$, so both of its set
maps are constant and its transformation profile is already in
the minimal band.  For a difference $p-q$, the two row operators
and the two column operators respectively agree.  If the common
row operator is nonzero, the common column operator is also
nonzero: $W$ is faithful over $A/I$ and $V$ factors through that
quotient.  A finite set map is uniquely determined by its nonzero
augmentation operator, so both maps agree literally.

If both operators are zero, both profiles belong to the band.
The remaining case has zero common row operator and nonzero
common column operator.  The certificate then supplies a literal
basis element $z$ with
\[
 pz=qz\ne0
\]
in the contracted monoid table.  In the original monoid the
products therefore equal the same outside element.  Since the
row maps of $p,q$ are constant, their composites on the right
with the row map of $z$ retain their respective constant values.
The displayed equality forces those values to agree.  The column
maps already agree because their augmentation operators are
nonzero.  Thus the profiles are equal in this last case too.
This step does not infer an equality of distinct constant maps
from equality of their zero augmentation operators.

The output
\path{code/outputs/d_six_square_zero_ideals.json} records all
ideal bases, quotient coordinates, spanning relations and actual
nonzero-product witnesses.  It proves the asserted spanning
equalities in $628$ cases.  Precisely eight ideals have no such
spanning list; they are retained below, not discarded.
The independent checker
\path{code/b_verify_six_square_zero_ideals.py} reconstructs all
$636$ corners and verifies their annihilation of $V$, their
two-sided ideal identities, and every literal spanning and
product relation.  It reproduces the same $628$ successes and
eight exceptions: the successful witnesses consist of $317$
basis-zero relations, $208$ differences with both operators
nonzero, and $103$ nonzero right-product repairs.  Its output is
\path{code/outputs/b_six_square_zero_ideals_verified.json}.

In a successful case the contracted algebra of the faithful
transformation-pair image is a quotient of $A/I$, hence hereditary
by Lemma~\ref{B:six-squarezero-target-modules}.  Its fourth rational
homology is zero by the band--Tor formula.  The map from the
original monoid to its transformation-pair image induces an
isomorphism in rational degree-two homology, since the two
balanced action presentations coincide.  Every degree-two class
pulls back from the image, so all its products vanish.
\end{proof}

\begin{proposition}[The eight remaining affine triangle families]
\label{B:six-squarezero-triangle-affine}
None of the eight remaining faithful triangle cases in
Lemma~\ref{B:six-squarezero-literal-ideals} can give the cohomology
ring of $S^2\times S^2$.
\end{proposition}
\begin{proof}
We describe the complete finite enumeration and its certificates.
Source indices refer to
\path{code/outputs/d_six_quotient_algebras.jsonl.gz}; the associated
ordinary zero-seven census indices are included to make the
identification independent of record order.  The coordinates of
the arrow ideal use the six distinguished nonzero basis elements
$q_0=1,q_1,\ldots,q_5$.
\[
\begin{array}{c|c|c|c}
\text{source}&\text{zero-seven class/opposite}&I&
 (\dim W_1,\dim W_2,\dim W_3)\\\hline
763&423/\text{no}&q_2-q_3-q_4+q_5&(1,1,3)\\
764&423/\text{yes}&q_2-q_3-q_4+q_5&(2,1,2)\\
1261&678/\text{no}&q_1-q_2+q_3&(1,1,3)\\
1262&678/\text{yes}&q_1-q_2+q_3&(2,1,2)\\
4187&2246/\text{no}&q_1-q_4+q_5&(2,1,2)\\
4188&2246/\text{yes}&q_1-q_4+q_5&(1,1,3)\\
4513&2415/\text{no}&q_1-q_4+q_5&(2,1,2)\\
4514&2415/\text{yes}&q_1-q_4+q_5&(1,1,3)
\end{array}
\]
The vertex order in the last column is the intrinsic triangle
order of Lemma~\ref{B:six-squarezero-target-modules}, rather than
the arbitrary order of the lifted idempotents in a data file.

The faithful-module statement is essential to the scope of this
enumeration.  None of the six distinguished basis vectors belongs
to $I$, as the quotient coordinates in the ideal certificate
verify.  Since $W$ is faithful over $A/I$, every outside basis
element consequently has a nonzero column augmentation operator:
its six-point column map is nonconstant.  Cases not satisfying this
faithfulness condition have already been excluded in the module
lemma by rank-one forms, or do not have the required Tor dimensions.
More explicitly, for $A/(a)$ the vertex dimensions are $(1,1,3)$:
if $b=0$, the tensor complex has degree-zero dimension three and
degree-one dimension one; if $c=0$, all balanced forms have rank
at most one.  For $A/(b)$ the dimensions are $(2,1,2)$:
if $a=0$, the degree-one and degree-two Tor dimensions are one
and two, respectively; if $c=0$, all balanced forms again have
rank at most one.  The only ways to lose faithfulness are these
vanishings, since all vertex spaces are nonzero and the two
remaining arrow actions occupy different Peirce blocks.

The column search \path{code/d_triangle_column_csp.cpp} starts
with all $6^6=46656$ functions on six points.  For each outside
basis element $q_i$ it retains the nonconstant functions with
the required number of fixed points
\[
 1+\operatorname{tr}(q_i\mid W)
 =1+\sum_j(\dim W_j)\chi_j(q_i),
\]
and all of its literal power relations.  The identity has its
forced map.  Every nonzero source product imposes the literal
composition equality for column maps, in their contravariant
order; every zero source product requires a constant composite.
Finally the displayed arrow relation is imposed exactly on the
five-dimensional augmentation matrices.

There is one symmetry normalization.  Choose an idempotent source
element and retain one representative of each conjugacy class of
its possible six-point idempotent maps.  These classes are exactly
the partitions of six given by the sizes of the fibers over its
fixed points.  The search branches over every remaining value;
each propagation step merely removes a value having no compatible
value in an explicitly imposed equation.  Thus this normalization
and propagation lose no simultaneous action.  The complete output
sets, before further stabilizer reduction, have sizes shown below.
\[
\begin{array}{c|r|r|r|r}
\text{source}&\text{column actions}&\text{column orbits}&
\text{row lifts}&\text{faithful full cases}\\\hline
763&816&39&1&26\\
764&300&62&2&56\\
1261&240&20&1&14\\
1262&85&27&1&10\\
4187&248&32&1&15\\
4188&96&20&1&14\\
4513&624&78&2&156\\
4514&288&52&3&156
\end{array}
\]
The second orbit reduction uses only the stabilizer of the fixed
idempotent map.  For rows, every nonzero augmentation operator
determines its three-point map uniquely.  For each zero operator
on an outside element the program tries all three constant maps
and retains exactly those satisfying all outside product equations.
The last column includes every resulting row lift.

The independent program
\path{code/a_triangle_column_audit.cpp}, with runner
\path{code/a_triangle_column_audit.py}, uses a different
multiplication-closure search.  It reconstructs the pools and
fiber-partition normalization and reproduces all eight complete
solution sets exactly, totaling $2697$ actions.  Its verified
output is \path{code/outputs/a_triangle_column_audit.json}.
Two further independent linkage checks,
\path{code/a_triangle24_coverage.py} and
\path{code/b_verify_triangle24_coverage.py}, verify all $330$
column orbits against these $2697$ actions, every constant row
lift, and the faithful-module rank filter.  They account for all
$93$ discarded column orbits and match the $237$ retained orbits
and their row lifts to exactly $447$ source certificates.  Their
outputs are \path{code/outputs/a_triangle24_coverage.json} and
\path{code/outputs/b_triangle24_coverage_verified.json}.

The cup calculation \path{code/d_triangle24_cups.py} first
retains the faithful $A/I$ branch by checking that the span of
its augmentation matrices has dimension five.  Low-rank tuples
are merely discarded as outside this branch; no new cup claim
is inferred from their rank.  The retained actions form precisely
the $447$ cases in the last column.  For each pair of actions the
full twenty-four-element monoid table is explicit: outside
products use the source table when nonzero; when that source
product is zero, the constant composite row and column maps
specify the unique band element.  Products involving the band
are the two actions or rectangular-band multiplication.  All
associativity equations are checked.
These are the full source tables, with six separate outside
labels and eighteen separate band labels.  Outside labels are
retained even when their transformation profiles agree.
Consequently the $447$ certificates cover nonfaithful source
actions as well; they are not tables obtained by replacing a
source with its smaller faithful image.  No finite-homotopy
assertion about such an image is required anywhere in this
enumeration.

For each monoid the program constructs an exact integral
projective resolution prefix through degree four, two integral
balanced forms forming a rational basis, and their affine normalized
two-cocycles.  It gives a literal module-linear comparison from
the projective prefix to the normalized bar resolution.  For all
four ordered products of the two degree-two classes, evaluation
on the image of \emph{every} degree-four projective generator is
zero.  Thus all four product cochains are identically zero on
this resolution; the conclusion is stronger than testing a few
homology cycles.  All comparisons, tables and cocycles are saved
in the $447$ compressed certificates named in
\path{code/outputs/d_triangle24_cups.json}.

The standalone checker \path{code/b_verify_triangle24_cups.py}
imports neither the resolution producer nor its bar-comparison
routine.  It verifies every table, exactness and integral
saturation of the projective prefix, every comparison equation
and its module translates, the full rational degree-two basis
by its pairing with tensor cycles, and products on a saturated
degree-four cycle basis.  All $447$ certificates pass, with the
same counts as the table, as recorded in
\path{code/outputs/b_triangle24_cups_verified.json}.
No conclusion about an uncomputed resolution tail is used.
Every faithful case has zero degree-two products, proving the
proposition.
\end{proof}

\begin{corollary}\label{B:six-squarezero-product-exclusion}
The $458$ square-zero row-module exceptions in
Proposition~\ref{B:six-row-census} cannot occur for a
twenty-four-element monoid model of $S^2\times S^2$.
\end{corollary}
\begin{proof}
Lemma~\ref{B:six-squarezero-target-modules} either makes all
degree-two forms rank one, in which case
Lemma~\ref{B:three-row-rank-one-products} kills every product,
or gives one of the faithful arrow quotients.  The literal ideal
gate and the eight retained affine families exhaust those
quotients by Lemma~\ref{B:six-squarezero-literal-ideals} and
Proposition~\ref{B:six-squarezero-triangle-affine}.
\end{proof}
